\documentclass[11pt,a4paper]{article}
\usepackage[margin=2.2cm]{geometry}
\usepackage{fontspec}
\setmainfont{Latin Modern Roman}
\setmonofont{Latin Modern Mono}
\usepackage[provide=*]{babel}
\babelprovide[import,main]{romanian}
\usepackage{xcolor}
\usepackage{booktabs,tabularx,array}
\usepackage{graphicx}
\usepackage{enumitem}
\usepackage{listings}
\usepackage[most]{tcolorbox}
\usepackage[hidelinks]{hyperref}
\usepackage{titlesec}
\usepackage{fancyhdr}
\usepackage{needspace}

\definecolor{accent}{HTML}{1F5F8B}
\definecolor{trap}{HTML}{B3261E}

\titleformat{\section}{\Large\bfseries\color{accent}}{}{0pt}{}[\vspace{0.2em}\color{accent}\titlerule]
\titlespacing*{\section}{0pt}{1.6em}{0.8em}
\titleformat{\subsection}{\large\bfseries\color{accent}}{}{0pt}{}
\titlespacing*{\subsection}{0pt}{1.1em}{0.4em}
\setlist{itemsep=0.25em,topsep=0.3em}
\setlength{\parindent}{0pt}
\setlength{\parskip}{0.4em}

\pagestyle{fancy}
\fancyhf{}
\renewcommand{\headrulewidth}{0.3pt}
\fancyhead[L]{Administrarea rețelelor}
\fancyhead[R]{\cursnr}
\fancyfoot[C]{\thepage}

\lstset{basicstyle=\ttfamily\small,backgroundcolor=\color{black!4},frame=single,
  rulecolor=\color{black!20},xleftmargin=0.5em,xrightmargin=0.5em,
  columns=fullflexible,keepspaces=true,showstringspaces=false,breaklines=true}

\newtcolorbox{capcana}[1][Capcane de examen]{enhanced,colback=trap!5,colframe=trap,
  fonttitle=\bfseries,title={\,#1},boxrule=0.6pt,arc=2pt,left=6pt,right=6pt}

\newcommand{\img}[3][0.85]{\IfFileExists{#2}{\begin{center}
  \includegraphics[width=#1\linewidth]{#2}\\[-0.2em]{\small\itshape #3}\end{center}}{}}
\newcommand{\cod}[1]{\texttt{#1}}
\newcolumntype{Y}{>{\raggedright\arraybackslash}X}
\newcommand{\titlu}[2]{\begin{center}{\Huge\bfseries\color{accent} #1}\\[0.4em]{\large #2}\end{center}}

% titlurile nu rămân singure la capăt de pagină
\let\oldsection\section
\renewcommand{\section}{\Needspace{9\baselineskip}\oldsection}
\let\oldsubsection\subsection
\renewcommand{\subsection}{\Needspace{6\baselineskip}\oldsubsection}

\newcommand{\cursnr}{Curs 4 -- DHCP}
\begin{document}
\titlu{DHCP}{Dynamic Host Configuration Protocol \\ Administrarea rețelelor \textbullet{} Curs 4}

\section{DHCP}
\begin{itemize}
  \item Protocol la \textbf{nivelul 7 (aplicație)}.
  \item Are rol de \textbf{asignare automată de adrese IP} către clienții din rețea.
  \item \textbf{2 componente:}
  \begin{itemize}
    \item \textbf{Server DHCP}, pe un calculator cu rol de server/router, configurat cu \textbf{IP fix, manual}.
    \item \textbf{Client DHCP}, pe celelalte calculatoare din aceeași rețea, \textbf{fără IP configurat}.
  \end{itemize}
  \item Serverul ține evidența adreselor alocate (\textbf{IP $+$ MAC}), ca să nu dea duplicate. Gestionarea se face \textbf{pe bază de MAC}.
  \item \textbf{LEASE TIME:} adresele sunt „închiriate” pe o perioadă. După expirare, clientul cere din nou, de obicei primește aceeași adresă dacă e liberă.
  \item \textbf{Transport:}
  \begin{itemize}
    \item Serverul ascultă pe \textbf{port 67}.
    \item Clientul pe \textbf{port 68}.
  \end{itemize}
  \item \textbf{Comunicația o inițiază clientul.}
  \item Folosește \textbf{UDP}.
  \item Pe lângă IP trimite și:
  \begin{itemize}
    \item \textbf{netmask (obligatorie)}
    \item gateway
    \item servere DNS
    \item servicii NTP
    \item BOOTP
    \item fișiere de boot
  \end{itemize}
  \item \textbf{Adrese dinamice:} dintr-un interval configurabil, se pot schimba la expirarea lease-ului sau când decide serverul, deci clientul poate primi altă adresă.
  \item \textbf{Rezervare:} pe baza \textbf{MAC-ului}, clientul primește mereu aceeași adresă. Se recomandă ca adresele rezervate să \textbf{nu facă parte din intervalul dinamic}.
\end{itemize}

\Needspace{28\baselineskip}
\section{Mesaje}
\renewcommand{\arraystretch}{1.35}
\small
\begin{tabularx}{\linewidth}{@{}p{2.9cm} p{1.3cm} Y@{}}
\toprule
\textbf{Mesaj} & \textbf{Cine} & \textbf{Detalii din slide}\\
\midrule
\textbf{DHCPDISCOVER} & client & începe procesul. \textbf{Broadcast L3 (255.255.255.255) și L2 (FF:..:FF)}, UDP \textbf{port 67}. Broadcast, pentru că \textbf{nu are IP și nu știe adresa serverului}\\
\textbf{DHCPOFFER} & server & ofertă cu un IP posibil. UDP \textbf{port 68}. L3: \textbf{broadcast} (clientul nu are IP). L2: \textbf{broadcast sau unicast} pe MAC-ul clientului, deja cunoscut\\
\textbf{DHCPREQUEST} & client & acceptă oferta. De obicei \textbf{broadcast}, ca să informeze pe ceilalți clienți și pe eventualele alte servere DHCP. Poate conține IP-ul din OFFER, \textbf{alt IP} sau \textbf{IP-ul vechi}\\
\textbf{DHCPACK} & server & confirmă REQUEST-ul. \textbf{Broadcast sau unicast}\\
\textbf{DHCPNACK} & server & \textbf{respinge} REQUEST-ul (IP invalid sau indisponibil). De obicei \textbf{broadcast}\\
\textbf{DHCPRELEASE} & client & anunță \textbf{eliberarea} adresei. \textbf{Unicast} către IP-ul serverului\\
\textbf{DHCPINFORM} & client & clientul \textbf{are deja IP} și cere setări adiționale. \textbf{Broadcast} dacă nu știe IP-ul serverului, \textbf{unicast} dacă îl știe. Serverul răspunde cu \textbf{DHCPACK} cu informațiile cerute\\
\textbf{DHCPDECLINE} & client & \textbf{refuză oferta} din OFFER, când detectează că IP-ul e \textbf{deja utilizat} în rețea\\
\bottomrule
\end{tabularx}
\normalsize

\bigskip
\begin{capcana}
\begin{itemize}
  \item \textbf{Server 67, client 68.} DISCOVER merge \textbf{la} 67, OFFER merge \textbf{la} 68.
  \item \textbf{DISCOVER e broadcast} pentru că clientul nu are IP și nu cunoaște serverul.
  \item \textbf{OFFER:} L3 broadcast, dar \textbf{L2 poate fi broadcast sau unicast}.
  \item \textbf{NACK $=$ server respinge. DECLINE $=$ client refuză.}
  \item \textbf{RELEASE $=$ unicast} către server. \textbf{REQUEST $=$ de obicei broadcast.}
  \item \textbf{INFORM:} clientul are deja IP, iar serverul răspunde cu \textbf{ACK}, nu cu OFFER.
  \item \textbf{Netmask e obligatorie}, restul (gateway, DNS) sunt opționale.
  \item \textbf{DHCP e L7}, nu L3 și nu L2.
  \item Rezervarea e pe \textbf{MAC}, nu pe IP.
\end{itemize}
\end{capcana}
\end{document}
